\chapter{Glossary and Sources}
\label{app:glossary}

\noindent
A bilingual glossary, and an account of where the material in this book came
from and how to check whether it still holds.

\chapterstub{
\textbf{Contents.} (E.1) Glossary, English to Polish --- checkpointer, durable
execution, human-in-the-loop, interrupt, resume, reducer, fan-out, fan-in, tool
calling, guardrails, grounding, trajectory, handoff, backpressure, superstep,
structured concurrency, context engineering, prompt caching. Included because the
documentation is in English and a large number of the conversations about it are
not. (E.2) Primary sources, with what each is good for and what it is stale
about. (E.3) Repositories worth reading rather than searching. (E.4) How to date
a piece of LangChain material in ten seconds: if it contains
\api{AgentExecutor}, \api{initialize\_agent}, \api{LLMChain} or
\code{langgraph.prebuilt.create\_react\_agent}, it predates October 2025 and
describes an architecture that no longer exists. (E.5) How this book was verified,
and where it was not --- listed explicitly, so a reader knows which parts to
distrust first.
}
